\section{Contesto e motivazioni}
\label{sec:contesto_motivazioni}


Negli ultimi anni si è assistito a una progressiva trasformazione
dell'architettura dei sistemi distribuiti: servizi e applicazioni,
tradizionalmente eseguiti in data center centralizzati, vengono
sempre più frequentemente spostati verso il bordo della rete.
Questa tendenza, nota come \emph{edge computing}, nasce dalla
necessità di ridurre la latenza percepita dagli utenti, supportare
applicazioni real-time e limitare il traffico verso il cloud
\cite{Shi2016_EdgeVisionChallenges}.
Il paradigma emergente non è più quello di un cloud monolitico,
ma di un \emph{continuum Edge--Cloud}, in cui risorse computazionali
eterogenee e geograficamente distribuite collaborano per soddisfare
i requisiti applicativi. In tale contesto, l'elaborazione può essere
collocata dinamicamente lungo la catena edge--cloud in funzione delle
condizioni di rete, della disponibilità di risorse e dei vincoli
applicativi.

\subsection*{Serverless Computing e FaaS nel continuum Edge--Cloud}

Parallelamente all'evoluzione verso l'edge, il paradigma
\emph{Serverless Computing} ha acquisito crescente rilevanza
come modello di programmazione cloud-native.
Nel modello serverless, lo sviluppatore distribuisce codice
senza occuparsi della gestione dell'infrastruttura sottostante,
che viene completamente delegata al provider
\cite{VanEyk2018_ServerlessIsMore}.
Il modello di fatturazione \emph{pay-as-you-go} e la scalabilità
automatica hanno contribuito alla sua ampia adozione nell'industria.

All'interno dell'ecosistema serverless, il modello
\emph{Function-as-a-Service} (FaaS) rappresenta la forma
più diffusa di esecuzione event-driven.
In una piattaforma FaaS, l'unità fondamentale di
esecuzione è una funzione stateless, invocata in risposta a
eventi o richieste e tipicamente eseguita all'interno di
container o ambienti isolati equivalenti.

L'estensione del paradigma FaaS al continuum Edge--Cloud
risulta particolarmente promettente: eseguire funzioni
direttamente su nodi edge consente di ridurre la latenza
di risposta e migliorare la qualità del servizio percepita.
Recenti lavori hanno mostrato come architetture FaaS
decentralizzate possano abilitare l'esecuzione distribuita
di funzioni tra edge e cloud, introducendo meccanismi di
offloading verticale e orizzontale \cite{Russo2023_Serverledge}.

\subsection{Le criticità energetiche dell'esecuzione all'edge}

Nonostante i benefici in termini di latenza e prossimità
ai dati, l'esecuzione all'edge introduce criticità significative.
I nodi edge sono tipicamente caratterizzati da risorse
computazionali limitate, eterogeneità hardware e condizioni
operative variabili. A differenza dei data center cloud,
progettati per offrire elevata capacità e ridondanza, i nodi
edge devono operare con vincoli stringenti su CPU, memoria e,
soprattutto, energia disponibile.

In molti scenari applicativi -- dispositivi alimentati a batteria,
infrastrutture soggette a limiti di potenza, contesti IoT -- l'energia
non rappresenta soltanto un costo operativo, ma una risorsa limitata
e dinamica. La letteratura sul resource scheduling in ambienti edge
evidenzia come la gestione efficiente delle risorse, inclusa
l'energia, costituisca una delle principali sfide ancora aperte
\cite{Luo2021_EdgeSchedulingSurvey,Kar2023_OffloadingFederatedSurvey}.

In tali contesti non è sufficiente adottare politiche che
minimizzano il consumo medio di energia: diventa necessario
controllare l'energia associata alle singole esecuzioni,
garantendo il rispetto di budget energetici prefissati
e potenzialmente variabili nel tempo.

\subsection{Il problema della granularità nella misurazione energetica}

Un prerequisito fondamentale per qualsiasi meccanismo di controllo
energetico è la disponibilità di misure affidabili e aggiornate a
granularità sufficientemente fine. Storicamente, la stima del consumo
energetico nei sistemi virtualizzati è stata affrontata a livello
di macchina fisica o di macchina virtuale, rendendo difficile
attribuire il costo energetico a singoli componenti applicativi.

Nel contesto FaaS, questa difficoltà è amplificata dalla natura
\emph{fine-grained} delle invocazioni e dalla variabilità delle
condizioni runtime, quali il carico concorrente, le interferenze
tra container e la distinzione tra esecuzioni a freddo
(\emph{cold start}) e a caldo (\emph{warm start}).
Negli ultimi anni, strumenti di osservabilità energetica per ambienti
containerizzati hanno reso disponibile una misurazione a granularità
significativamente più fine.

In particolare, Kepler (\emph{Kubernetes-based Efficient Power Level
Exporter}) stima il consumo energetico a livello di singolo processo
e container in ambienti Kubernetes, esportando metriche accessibili
in tempo reale tramite Prometheus \cite{Amaral2023_Kepler}.
La disponibilità di misurazioni energetiche a livello di container
costituisce un abilitatore fondamentale per integrare l'energia tra
le metriche decisionali delle piattaforme FaaS, aprendo la strada
a meccanismi di scheduling energeticamente consapevoli.

\subsection{Approximate Computing come leva per la riduzione dei consumi}

Una direzione promettente per la riduzione del consumo energetico
in ambito FaaS è rappresentata dal paradigma dell'\emph{Approximate
Computing}, che consente di ridurre il costo computazionale
accettando, in modo controllato, una diminuzione dell'accuratezza
del risultato prodotto. Lavori come JouleGuard~\cite{jouleguard}
hanno mostrato come sia possibile progettare meccanismi di controllo
runtime capaci di rispettare un budget energetico prefissato,
regolando dinamicamente il livello di approssimazione dell'esecuzione.

Applicata al paradigma FaaS, questa idea si traduce nella possibilità
di associare a una stessa funzione logica più varianti implementative,
caratterizzate da diversi profili energetici e da diversi livelli
di accuratezza del risultato. La selezione della variante da eseguire
può quindi diventare una decisione dello scheduler, guidata da
informazioni energetiche aggiornate a runtime e dai vincoli
espressi dall'utente al momento dell'invocazione.

\subsection{Verso uno scheduling energy-aware con consenso esplicito}

La combinazione tra (i)~l'adozione crescente di architetture edge--cloud
ibride, (ii)~i vincoli energetici stringenti nei nodi edge e
(iii)~la disponibilità di misurazioni energetiche a granularità
container-level rende oggi possibile e necessario studiare meccanismi
di esecuzione in grado di controllare dinamicamente l'energia
associata alle funzioni serverless.

Tuttavia, la riduzione del consumo energetico ottenuta tramite
approssimazione introduce un compromesso con l'accuratezza del
risultato che non può essere gestito unilateralmente dalla piattaforma.
Tale compromesso richiede un \emph{consenso esplicito} da parte
dell'utente che invoca la funzione: solo in presenza di tale consenso
lo scheduler è autorizzato a selezionare varianti approssimate.
Questa scelta progettuale garantisce che il comportamento del sistema
resti trasparente e controllabile, distinguendo nettamente tra
ottimizzazione energetica e degradazione indesiderata della qualità.

